\clearpage
\section*{The permit event study: comparison group and alderman turnover after 2015}
\textit{September 27, 2026. Agent-drafted checks requested by Jacob. The original-ward comparison was added to
Appendix Table C.2 at Jacob's request; the turnover checks are exploratory and change nothing in the paper.}

In the event study, each reassigned block is compared with every unchanged block within 500 feet of the old
boundary between its origin and destination wards, on both sides of that boundary: the ward-pair-by-year fixed
effect pools them. Unchanged blocks in the destination ward had the new alderman before and after 2015, so they are
valid controls only if trends are parallel across the old line. Comparing reassigned blocks only with unchanged
blocks in the ward they left (ward-pair-by-original-ward-by-year fixed effects) gives $-0.132$ (0.066) for the
combined effect, $-0.261$ (0.126) for moves toward greater stringency and $0.049$ (0.073) for moves toward greater
leniency, against $-0.127$, $-0.231$ and $0.067$ in the main specification, with larger standard errors.

A larger issue is timing. Treatment is fixed at the aldermen serving the two wards in 2014, and the sample requires
only that both still served in June 2015; blocks are then followed through 2020. After the May 2019 election, 67
percent of reassigned blocks and 55 percent of unchanged blocks have a different alderman in at least one ward of
their pair, against 3 to 6 percent in 2016--2018. The main estimate is concentrated in these later years: $-0.087$
(0.064) in 2015--2018 and $-0.225$ (0.063) in 2019--2020. Ending the window in 2018, or dropping block-years from a
pair's first change, gives $-0.088$ and $-0.083$, neither significant.

Classifying blocks each year by the aldermen actually serving requires the 2006--2022 index, because aldermen first
elected in 2019 have no 2006--2014 score. On that index, the fixed 2014 classification gives $-0.100$ (0.054) pooled,
$-0.034$ in 2015--2018 and $-0.263$ in 2019--2020; the classification by serving aldermen gives $0.029$ (0.062) pooled
and essentially zero in both periods. The 2019 turnover reversed the direction of 203 of the 687 reassigned blocks.
Under the fixed classification, the 2019--2020 effect is $-0.169$ (0.071) where the direction held and $-0.485$
(0.081) where it reversed. A block given a more lenient alderman in 2019 should gain permits if serving aldermen's
stringency drives permitting; these blocks show the largest decline relative to their 2014 classification.

The permit decline therefore follows the 2015 assignment, not the stringency of the aldermen serving later. Possible
explanations are persistent deterrence of projects from 2015--2018, other changes in the reassigned areas after
2019, or a small number of boundaries; the reversed-direction groups contain 97 and 106 blocks. Which boundaries
carry the 2019--2020 effect has not yet been examined. The paper's description of this result, and the stable-sample
definition, should be revisited before the manuscript is rewritten.

\textit{Reproduction.} The original-ward comparison is built by
\path{tasks/run_event_study_permit/code/build_binary_event_study_table.R}. The turnover checks run with
\texttt{make} in \path{tasks/audits/permit_event_study_turnover/code/}, which reads the production block panel, the
monthly alderman panel and the 2006--2022 index; the main estimate is reproduced before any change. A ward-month
vacancy counts as a change; the serving alderman of a ward-year is the one serving most of its months.
